\section{Limitations and threats to validity}
\label{sec:limitations}

A machine-checked proof removes one risk, an error in the argument. It
leaves the others in place. We list them from the most to the least
important.

\paragraph{The Dex side is a model, not the code.}
The semantics is written by hand from the server's source and validated by
review. It is not extracted from the implementation, and no tool checks
that the server refines it. A mistake in reading the code is a mistake in
the theorem's meaning. Three things reduce the risk and none removes it.
Every rule cites the implementing code. Three adversarial reviews checked
the rules against that code. And the model's executable form runs the
examples and the lockstep checks, although only against itself, not against
a running server. The Lean model is, further, a \emph{restriction} of the
specification to what the translation and the admissible environment use;
each omission is listed in the development with the reason why it cannot
matter. The specification also pins a base commit of the server. The server
has changed since, including the two fixes of
\S\ref{sec:server-bugs}; the second changes what happens to the messages of
a canceled Step, a behavior that a translated Flow reaches only at its
close, where it is unobservable. That last sentence is an argument, not a
theorem.

\paragraph{Failure-free environment.}
There are no crashes, no Worker failures, no retries, no timeouts, and no
redelivery in the model. Durability is the reason a workflow engine exists,
and this theorem says nothing about it. The intended complement is a
separate proposition, correctness under crashes: roughly, that crashes,
retries, and redelivery add no observable behavior beyond what the
failure-free semantics allows, apart from repeated external effects. Its
precise statement is itself future work. The server has a
deterministic simulation harness with injected crashes, which is relevant
evidence but is not a proof.

\paragraph{Fairness is assumed, not provided.}
Liveness holds for fair runs, and the fairness conditions of
Table~\ref{tab:fairness} are assumptions about components outside the
rules. F2 assumes that the timer service eventually fires a due timer. F3
assumes that the dispatcher eventually starts a pending method. F5 assumes
that the engine eventually commits or drops a returned response. F4 is the
one most worth stating plainly: it assumes that every Worker task
eventually settles, by returning or by being abandoned after its invocation
was canceled. The server does not guarantee this on its own. Cancellation
of a running handler is cooperative, and in the failure-free model there is
no method timeout to end a handler that never returns. A deployment in
which a handler can hang is outside the liveness theorem; the safety
theorem still applies to it. The fair runs that the development exhibits
(\S\ref{sec:nonvacuity}) cover termination, deadlock, and divergence by an
infinite input word. A fair run of a program that diverges on its own, on a
finite word, is not exhibited, and no lemma says that every reachable
configuration has a fair continuation.

\paragraph{A restricted environment.}
The admissible environment starts one Flow and publishes plain data to
external Channels. It makes no RPC, never stops a Flow, never sets an
Attribute from outside, and never deletes a message or skips a timer. The
specification gives rules for all of these, but they are outside the
theorem. In particular an RPC can observe intermediate state and can fail
with a lock conflict, and stopping a Flow cancels Steps at arbitrary
points; both would need \Lang{} counterparts.

\paragraph{Well-formed programs only.}
The hypothesis \lean{Program.IsWellFormed} excludes programs that can break
a mutex requirement. For such a program the two sides differ: \Lang{}
refuses the step, and the Flow takes it. The scope is semantic and
undecidable in general. The static lock discipline of
\S\ref{sec:nonvacuity} is sufficient, not necessary. It tracks statically
known mutexes only, and it empties the lock set after a call and after a
\kw{try}. So it rejects \kw{unlock}, \kw{wait}, \kw{signal}, and
\kw{broadcast} on a created mutex, and a program that holds a lock across a
call and unlocks after it, although such programs can be well formed. For
them well-formedness must be proved by other means.

\paragraph{The timing model.}
\Lang{} counts stamped actions exactly as the translation counts commits
(\S\ref{sec:timing}). The theorem is exact for this model and silent about
any other. A program whose observable behavior depends on the order of two
deadlines that are close together can behave differently under another cost
model. We consider such programs incorrect, but the theorem does not
exclude them; it just assigns them the Dex timing. The equivalence is also
untimed: it does not say that an event happens within some time, only that
the same sequences of events are possible. Durations are whole seconds.

\paragraph{Unbounded counters.}
Step numbers and timer numbers are unbounded in the model and 32-bit in the
server. The order of timers with equal fire times was derived from the
protobuf encoding of timer IDs by hand; that derivation is an assumption of
the model, and it holds for fewer than $2^{32}$ Step executions of one type
per run.

\paragraph{\Lang{} is a chosen language.}
\Lang{}'s semantics was designed so that the equivalence is exact: FIFO
mutexes and condition variables, cooperative cancellation at four
cancellation points, asynchronous \kw{emit}, a \kw{select} with one receive
branch and one timeout branch, late timers, and a single-use join. Each
choice is defensible and close to common practice, but \Lang{} is not an
existing language, and a different reasonable choice (say, a cancellable
\kw{lock}) would need a different translation and a new proof. The theorem
shows that Dex can express \emph{a} conventional concurrent language. It
does not show that every feature of every such language has a faithful
encoding.

\paragraph{Observations.}
Only inputs, external effects, and the program's result are observable.
Reads of a Flow's state by a client are not, and the events of the Flow's
internal Channels are projected away. An observer with access to those
could tell a translated Flow from its source program.

\paragraph{Trust in the statements.}
The kernel checks proofs, not intentions. A reader who wants to rely on the
theorem must read the model, the translation, and the statements, about
3,450 lines, against the specification. The development maps every section
and every numbered assumption of the specification to its Lean definition
to make that reading possible.
